indocumentclass[12pt,a4paper]{article}

% =====================================================
% Packages
% =====================================================

\usepackage[utf8]{inputenc}
\usepackage[T1]{fontenc}
\usepackage[brazilian]{babel}
\usepackage{times}
\usepackage{geometry}
\geometry{margin=2cm}

\usepackage{graphicx}
\usepackage{float}
\usepackage{caption}
\usepackage{hyperref}
\usepackage{enumitem}
\usepackage{fancyhdr}
\usepackage{xcolor}
\usepackage{listings}
\usepackage{longtable}
\usepackage{array}
\usepackage{booktabs}
\usepackage[numbers]{natbib}
\setcounter{secnumdepth}{5}
% =====================================================
% Code blocks
% =====================================================

% =====================================================
% Better terminal listings
% =====================================================

\usepackage{listings}
\usepackage{xcolor}
\usepackage{upquote}      % keeps ASCII ', " and ` when copying
\usepackage{inconsolata}  % high-quality monospace font

% Colors
\definecolor{cmdbg}{RGB}{20,20,20}
\definecolor{cmdtext}{RGB}{245,245,245}

\definecolor{outbg}{RGB}{245,245,245}
\definecolor{outtext}{RGB}{0,0,0}

\definecolor{linenumber}{RGB}{130,130,130}

% Common settings
\lstset{
    basicstyle=\ttfamily\small,
    upquote=true,
    breaklines=true,
    showstringspaces=false,
    keepspaces=true,
    columns=fullflexible,
    tabsize=4,
    numbers=left,
    numberstyle=\small\color{linenumber},
    numbersep=8pt,
    frame=single,
    rulecolor=\color{black},
    captionpos=b
}

% Command blocks
\lstdefinestyle{command}{
    language=bash,
    backgroundcolor=\color{cmdbg},
    basicstyle=\ttfamily\color{cmdtext},
    keywordstyle=\color{cmdtext},
    commentstyle=\color{green},
    stringstyle=\color{yellow},
    identifierstyle=\color{cmdtext},
    numberstyle=\color{linenumber},
}

% Terminal output blocks
\lstdefinestyle{output}{
    language={},
    backgroundcolor=\color{outbg},
    basicstyle=\ttfamily\small\color{outtext},
    keywordstyle=\color{outtext},
    commentstyle=\color{outtext},
    stringstyle=\color{outtext},
    identifierstyle=\color{outtext},
    numberstyle=\small\color{linenumber},
}
% =====================================================
% Header / Footer
% =====================================================

\pagestyle{fancy}
\fancyhf{}

\fancyhead[L]{Relatório Técnico --- <NOME DO CTF>}
\fancyhead[R]{<AUTORES>}
\fancyfoot[C]{\thepage}

\renewcommand{\headrulewidth}{0.4pt}

% =====================================================
\begin{document}

%%%%%%%%%%%%%%%%%%%%%%%%%%%%%%%%%%%%%%%%%%%%%%%%%%%%%%%%
%%%%%%%%%%%%%%%%%%%%%%%% CAPA %%%%%%%%%%%%%%%%%%%%%%%%%%%
%%%%%%%%%%%%%%%%%%%%%%%%%%%%%%%%%%%%%%%%%%%%%%%%%%%%%%%%

\begin{titlepage}

\centering

\vspace*{2cm}

{\Huge\bfseries <NOME DO CTF>\par}

\vspace{0.7cm}

{\Large Relatório Técnico}

\vspace{2cm}

\begin{tabular}{rl}
Disciplina:& <DISCIPLINA>\\
Professor:& <PROFESSOR>\\
Integrantes:& <NOMES>\\
Data:& \today
\end{tabular}

\vfill

\end{titlepage}

\tableofcontents

\newpage

%%%%%%%%%%%%%%%%%%%%%%%%%%%%%%%%%%%%%%%%%%%%%%%%%%%%%%%%
%%%%%%%%%%%%%%%%%%%% INTRODUÇÃO %%%%%%%%%%%%%%%%%%%%%%%%%
%%%%%%%%%%%%%%%%%%%%%%%%%%%%%%%%%%%%%%%%%%%%%%%%%%%%%%%%

\section{Introdução}

\subsection{Objetivo}

<DESCREVER O OBJETIVO DO DESAFIO>

\vspace{1em}

\subsection{Contextualização}

<DESCREVER O CENÁRIO DO CTF>

\vspace{1em}

\subsection{Tecnologias utilizadas}

\begin{itemize}
    \item Ubuntu
    \item Docker
    \item Apache
    \item FTP
    \item FileBrowser
    \item <OUTRAS TECNOLOGIAS>
\end{itemize}

%%%%%%%%%%%%%%%%%%%%%%%%%%%%%%%%%%%%%%%%%%%%%%%%%%%%%%%%
%%%%%%%%%%%%%%%% EXECUÇÃO DO AMBIENTE %%%%%%%%%%%%%%%%%%
%%%%%%%%%%%%%%%%%%%%%%%%%%%%%%%%%%%%%%%%%%%%%%%%%%%%%%%%

\section{Execução do Ambiente}

\subsection{Arquivos entregues}

\begin{itemize}
    \item Dockerfile
    \item docker-compose.yml
    \item Scripts auxiliares
    \item Configurações
    \item Arquivos do desafio
\end{itemize}

\subsection{Estrutura do projeto}

\begin{lstlisting}


\end{lstlisting}

\subsection{Execução}

\begin{lstlisting}[style=command]
docker compose build

docker compose up -d
\end{lstlisting}

\end{lstlisting}
\subsection{Serviços}

\begin{longtable}{ll}
\toprule
Porta & Serviço\\
\midrule
8080 & Apache\\
<PORTA> & FTP\\
<PORTA> & FileBrowser\\
\bottomrule
\end{longtable}

%%%%%%%%%%%%%%%%%%%%%%%%%%%%%%%%%%%%%%%%%%%%%%%%%%%%%%%%
%%%%%%%%%%%%%%%%%%%%% ARQUITETURA %%%%%%%%%%%%%%%%%%%%%%
%%%%%%%%%%%%%%%%%%%%%%%%%%%%%%%%%%%%%%%%%%%%%%%%%%%%%%%%

\section{Arquitetura do Cenário}

<INSERIR DIAGRAMA>

\begin{figure}[H]
\centering
\includegraphics[width=.9\linewidth]{architecture.png}
\caption{Arquitetura do ambiente}
\end{figure}

%%%%%%%%%%%%%%%%%%%%%%%%%%%%%%%%%%%%%%%%%%%%%%%%%%%%%%%%
%%%%%%%%%%%%%%%%%%%% VULNERABILIDADES %%%%%%%%%%%%%%%%%%
%%%%%%%%%%%%%%%%%%%%%%%%%%%%%%%%%%%%%%%%%%%%%%%%%%%%%%%%

\section{Vulnerabilidades Implementadas}

\subsection{Vulnerabilidade 1}

<TIPO>

Descrição:

<DESCREVER>

\vspace{1em}

\subsection{Vulnerabilidade 2}

<DESCREVER>

\vspace{1em}

\subsection{Escalada de privilégios}

<DESCREVER>

%%%%%%%%%%%%%%%%%%%%%%%%%%%%%%%%%%%%%%%%%%%%%%%%%%%%%%%%
%%%%%%%%%%%%%%%%%%%%%% WALKTHROUGH %%%%%%%%%%%%%%%%%%%%%
%%%%%%%%%%%%%%%%%%%%%%%%%%%%%%%%%%%%%%%%%%%%%%%%%%%%%%%%

\section{Walkthrough}

%%%%%%%%%%%%%%%%%%%%%%%%%%%%%%%%%%%%%%%%%%%%%%%%%%%%%%%%

\subsection{Reconhecimento}

\subsubsection{Nmap}
\begin{lstlisting}[style=command]
nmap -A -T5 localhost
\end{lstlisting}

Observações:

\begin{lstlisting}[style=output]
PORT      STATE SERVICE         VERSION
...
8080/tcp  open  http            Apache httpd 2.4.66 ((Ubuntu))
|_http-open-proxy: Proxy might be redirecting requests
|_http-server-header: Apache/2.4.66 (Ubuntu)
|_http-title: Begin the story
...
\end{lstlisting}

\subsubsection{Gobuster}
Como o servidor utiliza Apache, foi realizada enumeração de diretórios para arquivos `html,php,js`.

\begin{lstlisting}[style=command]
gobuster dir -u http://localhost:8080 -w Downloads/DirBuster-2007_directory-list-lowercase-2.3-small.txt -x html,php,js -t 1000
\end{lstlisting}
\begin{lstlisting}[style=output, label={lst:gobuster}]]
...
===============================================================
Starting gobuster in directory enumeration mode
===============================================================
/.html                (Status: 403) [Size: 316]
/index.html           (Status: 200) [Size: 1539]
/run.html             (Status: 200) [Size: 9765]
/.html                (Status: 403) [Size: 316]
Progress: 326572 / 326576 (100.00%)
...
\end{lstlisting}


\subsubsection{Porta 8080}
Em http://localhost:8080/
\paragraph{index.html}
A página inicial apresenta um vídeo
\begin{figure}[H]
    \centering
    \includegraphics[width=0.75\linewidth]{8080 index.png}
    \caption{Página inicial localhost:8080 }
    \label{fig:placeholder}
\end{figure}

Feito o download pelo HTML e explorado com exiftool:

\begin{lstlisting}[style=command]
exiftool welcome.mp4
\end{lstlisting}

% Observações:

\begin{lstlisting}[style=output]
...
Comment                         : /scene0
XMP Toolkit                     : Image::ExifTool 12.76
Artist                          : You wanna watch a drama? A Drama can be written from anyone's life, seek your own story
...
\end{lstlisting}

Esse vídeo parece ser a cena 0 de um 'Drama'.

\vspace{1em}

Também foi realizada uma requisição HTTP:

\begin{lstlisting}[style=command]
curl -I http://localhost:8080/
\end{lstlisting}

\begin{figure}[H]
\centering
\includegraphics[width=.6\linewidth]{curl localhost 8080 .png}
\caption{Resposta HTTP}
\end{figure}

Foi identificado o cabeçalho \texttt{X-Hint: /scene1}


\vspace{1em}
% Todo make all paragraphs hyperrefs like this
% todo change hyperref box color to a deep discrete dark blueish purple 
\paragraph{\href{http://localhost:8080/Scene1.html}{8080/Scene1.html}}

Na página:

\begin{figure}[H]
    \centering
    \includegraphics[width=0.75\linewidth]{scene 1.png}
    \caption{\texttt{localhost:8080/scene1.html}}
    \label{fig:placeholder}
\end{figure}


Foi feito download do vídeo, pelo HTML da página. 
Note que a pessoa é chamada de "aluna artista" no texto da página.
Analizando os metadados do vídeo:

\begin{lstlisting}[style=command]
exiftool strange_phenomenon_aeri.mp4
\end{lstlisting}
\begin{lstlisting}[style=output]
...
Encoder : $6$<yummy>$
XMP Toolkit :
Yp2586tEONHHNO8K2K4eAH6iEajQEV6HHPg65vkXQFWLGVO9K9fg.TWDrhqsqXlgLpV55qwAgOUGlEsO//EY//
Artist :
YWVzcGFGaWxlcw==
...
\end{lstlisting}

Decodificando o valor de 'Artist':

\begin{lstlisting}[style=command]
echo -n "YWVzcGFGaWxlcw==" | base64 -d
\end{lstlisting}
\begin{lstlisting}[style=output]
aespaFiles
\end{lstlisting}

Parecia existir um hash dividido entre metadados 'Encoder' e 'XMP Toolkit'. 

\vspace{1em}

Ao clicar na página aparece um GIF:
\begin{figure}[H]
\centering
\includegraphics[width=.6\linewidth]{scene1 gif .png}
\caption{GIF exibido}
\end{figure}

Inspecionando o HTML:

\begin{figure}[H]
\centering
\includegraphics[width=.6\linewidth]{see scene1 html calling gif.png}
\caption{HTML referenciando o GIF}
\end{figure}

Foi possível realizar download acessando http://localhost:8080/SALTy\_and\_sweet.gif:

Analisando com strings:

\begin{lstlisting}[style=command]
strings SALTy_and_sweet.gif
\end{lstlisting}
\begin{lstlisting}
...
;l'existence precede l'essence==
lexistence
...
\end{lstlisting}

Como SALT está em caixa alta, e tem um hash no vídeo, supõe-se que este contém SALT para o hash. Das duas frases inseridas, foi tentado primeiro o SALT menor.

Hash reconstruída:
% TODO FIX THIS INTRODUCES A NEW LINE IN HASH. CANNOT OVERFLOW BUT CANNOT INTRODUCE NEWLINE 
\begin{lstlisting}
$6$lexistence$Yp2586tEONHHNO8K2K4eAH6iEajQEV6HHPg65vkXQFWLGVO9K9fg.TWDrhqsqXlgLpV55qwAgOUGlEsO//EY//
\end{lstlisting}

\begin{lstlisting}[style=command]
hashid -m '<HASH>'
\end{lstlisting}

Observação: 'cooper' é uma entrada na 'rockyou.txt'

Wordlist:
\begin{lstlisting}[style=command]
echo cooper >> p.txt

hashcat -m 1800 -a 0 hash.txt p.txt
\end{lstlisting}
\begin{lstlisting}[style=output]
$6$lexistence$Yp2586tEONHHNO8K2K4eAH6iEajQEV6HHPg65vkXQFWLGVO9K9fg.TWDrhqsqXlgLpV55qwAgOUGlEsO//EY//:cooper
\end{lstlisting}

Assim:

\begin{itemize}
\item Usuário: aespaFiles
\item Senha: cooper
\end{itemize}

Estas credenciais são utilizadas posteriormente no FTP.

%%%%%%%%%%%%%%%%%%%%%%%%%%%%%%%%%%%%%%%%%%%%%%%%%%%%%%%%

\paragraph{\href{http://localhost:8080/Scene1.html}{8080/Run.html}} foi acessado seguindo o gobuster \ref{lst:gobuster}. Possui múltiplas imagens que podem ser movidas

\begin{figure}[H]
    \centering
    \includegraphics[width=0.5\linewidth]{8080 run 2 .png}
    \caption{8080/Run.html}
    \label{fig:placeholder}
\end{figure}

Todas fotos foram baixadas pelo HTML, e analizadas com comando \texttt{strings}.

Sabe-se pelo HTML como chegar nos links das imagens 
\begin{figure}[H]
    \centering
    \includegraphics[width=0.75\linewidth]{run html code where are the images.png}
    \caption{HTML com fonte da imagem em \texttt{src="WDA\_karina.jpeg"}}
    \label{fig:placeholder}
\end{figure}

Segue a resolução dessa parte:
\begin{enumerate}
\item Baixar todas as imagens.
\item Executar strings em cada imagem.
\item Decodificar a mensagem em Base64.
\item A Ningning (mulher de branco no cavalo, imagem escondida ao fundo) possui dois hashes.
\item Essas hashes serão utilizadas posteriormente como passphrases do steghide.
\begin{enumerate}
\item Primeira hash:
    \begin{itemize}
        \item SHA512
        \item Descoberta pela CrackStation
    \end{itemize}
\item Segunda hash:
    \begin{itemize}
        \item bcrypt
        \item descoberta por hashcat e rockyou
    \end{itemize}
\end{enumerate}
\end{enumerate}


\subparagraph{WDA\_karina.jpeg}, Imagem em http://localhost:8080/WDA\_karina.jpeg
\begin{lstlisting}[style=command]
 strings WDA_karina.jpeg 
\end{lstlisting}
\begin{lstlisting}[style=output]
...
_.ha<x1
9MK3
T3VyIGJlYXV0aWZ1bCBjcmVhdHVyZSBzZWFyY2hlcyB1bmRlciAxMDAwbQ==
\end{lstlisting}
\begin{lstlisting}[style=command]
echo "T3VyIGJlYXV0aWZ1bCBjcmVhdHVyZSBzZWFyY2hlcyB1bmRlciAxMDAwbQ==" | base64 -d
\end{lstlisting}
\begin{lstlisting}[style=output]
Our beautiful creature searches under 1000m
\end{lstlisting}

\begin{figure}[H]
    \centering
    \includegraphics[width=0.5\linewidth]{8080 run 1 .png}
    \caption{WDA\_karina.jpeg -- \textit{Our Beautiful Creature was flying under 1000m}}
    \label{fig:placeholder}
\end{figure}

\subparagraph{WDA\_winter.webm}, mesmo processo da karina

\subparagraph{WDA\_Ningning.jpeg},

\begin{lstlisting}[style=command]
strings WDA_Ningning.jpeg 
\end{lstlisting}
\begin{lstlisting}[style=output]
...
q,9G
Fq+b
PNWO
NDU3ODhmNjM3ZjUzZjRlYWUyNTZjYTNlYTQ0YTY5ZDg3MzlmNjUwMTdiYTc0ZjUyNWUzZGM3ZTVj
OWZlZjE4NjRjNWVjM2NmZTA3Nzg4NjYzZjY2NGQxOTM1ODdhNjU2MGRmOTA3N2I0NzljNTI5ODUw
ZDUwMzBlMjUyYjEwMzE=
JDJ5JDEwJFFlWk9IejI5UFI1OUtpakdjeFIvYnUwQmNxS3ZXbzIzSVk3bGR1eGJLaXdxSUo4YUJq
dHZxCg==
\end{lstlisting}
Parecem ter duas entradas em base64:
\begin{lstlisting}[style=command]
echo "NDU3ODhmNjM3ZjUzZjRlYWUyNTZjYTNlYTQ0YTY5ZDg3MzlmNjUwMTdiYTc0ZjUyNWUzZGM3ZTVj
OWZlZjE4NjRjNWVjM2NmZTA3Nzg4NjYzZjY2NGQxOTM1ODdhNjU2MGRmOTA3N2I0NzljNTI5ODUw
ZDUwMzBlMjUyYjEwMzE=" | base64 -d
\end{lstlisting}
%TODO DONT OVERFLOW PDF 
\begin{lstlisting}[style=output]
45788f637f53f4eae256ca3ea44a69d8739f65017ba74f525e3dc7e5c9fef1864c5ec3cfe07788663f664d193587a6560df9077b479c529850d5030e252b1031
\end{lstlisting}
% TODO ADD CRACKSTATION REF
Parece um hash. Jogando no crackstation
\begin{figure}
    \centering
    \includegraphics[width=0.75\linewidth]{crackstation .png}
    \caption{Hash quebrado -- \textit{complexity}}
    \label{fig:placeholder}
\end{figure}

Estudando a outra string:
%TODO DONT OVERFLOW PDF 
\begin{lstlisting}[style=command]
echo "JDJ5JDEwJFFlWk9IejI5UFI1OUtpakdjeFIvYnUwQmNxS3ZXbzIzSVk3bGR1eGJLaXdxSUo4YUJqdHZxCg==" | base64 -d
\end{lstlisting}
\begin{lstlisting}[style=output]
$2y$10$QeZOHz29PR59KijGcxR/bu0BcqKvWo23IY7lduxbKiwqIJ8aBjtvq
\end{lstlisting}

\subparagraph{Quebrando a hash com salt}, usando entrada da rockyou e hashcat
\begin{lstlisting}[style=command]
 hashid -m '$2y$10$QeZOHz29PR59KijGcxR/bu0BcqKvWo23IY7lduxbKiwqIJ8aBjtvq'
\end{lstlisting}
\begin{lstlisting}[style=output]
...
[+] bcrypt [Hashcat Mode: 3200]
\end{lstlisting}
Nota: 'supernovaflyingbakekang' está na rockyou.txt. Aqui foi criado uma wordlist apenas para demonstrar a quebra, espera-se que as atacantes usem rockyou.txt
\begin{lstlisting}[style=command]
echo "supernovaflyingbakekang" > wordlist.txt

echo '$2y$10$QeZOHz29PR59KijGcxR/bu0BcqKvWo23IY7lduxbKiwqIJ8aBjtvq' > hash.txt

hashcat -m 3200 -a 0 hash.txt wordlist.txt
\end{lstlisting}

Após a quebra, pode-se encontrar a senha novamente com 

\begin{lstlisting}[style=command]
hashcat -m 3200 -a 0 hash.txt wordlist.txt --show
\end{lstlisting}
\begin{lstlisting}[style=output]
$2y$10$QeZOHz29PR59KijGcxR/bu0BcqKvWo23IY7lduxbKiwqIJ8aBjtvq:supernovaflyingbakekang
\end{lstlisting}

%%%%%%%%%%%%%%%%%%%%%%%%%%%%%%%%%%%%%%%%%%%%%%%%%%%%%%%%

\subsection{FTP}

A imagem da karina disse sobre search under 1000. Fez-se nmap under ports 1000


Notas atuais:

\begin{lstlisting}
binwalk -e complaexity.mp4

steghide extract \
-sf image.jpg \
-p complexity
\end{lstlisting}

Resultado:

\begin{lstlisting}
who_are_they.txt
\end{lstlisting}

Conteúdo:

\begin{lstlisting}
user: complexity

password: complexitytrailer

Port of file-browser
\end{lstlisting}

%%%%%%%%%%%%%%%%%%%%%%%%%%%%%%%%%%%%%%%%%%%%%%%%%%%%%%%%

\subsection{FileBrowser}

<PLACEHOLDER>

%%%%%%%%%%%%%%%%%%%%%%%%%%%%%%%%%%%%%%%%%%%%%%%%%%%%%%%%

\subsection{Acesso Inicial}

<PLACEHOLDER>

%%%%%%%%%%%%%%%%%%%%%%%%%%%%%%%%%%%%%%%%%%%%%%%%%%%%%%%%

\subsection{Enumeração Local}

<PLACEHOLDER>

%%%%%%%%%%%%%%%%%%%%%%%%%%%%%%%%%%%%%%%%%%%%%%%%%%%%%%%%

\subsection{Escalada de Privilégios}

<PLACEHOLDER>

%%%%%%%%%%%%%%%%%%%%%%%%%%%%%%%%%%%%%%%%%%%%%%%%%%%%%%%%

\subsection{Flag Root}

<PLACEHOLDER>

%%%%%%%%%%%%%%%%%%%%%%%%%%%%%%%%%%%%%%%%%%%%%%%%%%%%%%%%
%%%%%%%%%%%%%%%%%%%% IMPLEMENTAÇÃO %%%%%%%%%%%%%%%%%%%%%
%%%%%%%%%%%%%%%%%%%%%%%%%%%%%%%%%%%%%%%%%%%%%%%%%%%%%%%%

\section{Detalhes de Implementação}

\subsection{Decisões de projeto}

<PLACEHOLDER>

\subsection{Desafios encontrados}

<PLACEHOLDER>

\subsection{Melhorias futuras}

<PLACEHOLDER>

%%%%%%%%%%%%%%%%%%%%%%%%%%%%%%%%%%%%%%%%%%%%%%%%%%%%%%%%
%%%%%%%%%%%%%%%%%%%% CONCLUSÃO %%%%%%%%%%%%%%%%%%%%%%%%%
%%%%%%%%%%%%%%%%%%%%%%%%%%%%%%%%%%%%%%%%%%%%%%%%%%%%%%%%

\section{Conclusão}

<PLACEHOLDER>

%%%%%%%%%%%%%%%%%%%%%%%%%%%%%%%%%%%%%%%%%%%%%%%%%%%%%%%%
%%%%%%%%%%%%%%%%%%%% APÊNDICES %%%%%%%%%%%%%%%%%%%%%%%%%
%%%%%%%%%%%%%%%%%%%%%%%%%%%%%%%%%%%%%%%%%%%%%%%%%%%%%%%%

\appendix

\section{Estrutura do Docker}

<PLACEHOLDER>

\section{Scripts}

<PLACEHOLDER>

\section{Arquivos auxiliares}

<PLACEHOLDER>

\bibliographystyle{unsrtnat}
\bibliography{bib}

\end{document}do the ctf 
```
